\documentclass[11pt]{article}
\usepackage[margin=1in]{geometry}
\usepackage{amsmath,amssymb,amsthm,mathtools}
\usepackage[hidelinks]{hyperref}
\newtheorem{theorem}{Theorem}
\newtheorem{lemma}[theorem]{Lemma}
\newtheorem{corollary}[theorem]{Corollary}
\newcommand{\Q}{\mathbb Q}
\newcommand{\C}{\mathbb C}
\newcommand{\F}{\mathbb F}
\title{Effective fixed-multiplier cutoffs and nearby-prime residual degree sets}
\author{User-supplied refinements checked and integrated; no priority claim}
\date{16 November 2026}
\begin{document}
\maketitle
Write $F_n(T)=\sum_{j=0}^n\binom nj(6j)!T^j$ with ordinary factorial
shape one. The \href{PRIME_GAPS_AND_SMALL_KERNELS.tex}{prime-gap and
small-kernel proof} retains the exposed primitive edge, the even
square-divisor bound and integer first-subresultant certificates for
offsets through twelve. This note makes its eventual statement effective
and retains full residual factor-degree information. General all-index
SFC, arbitrary Gamma shapes and independent/formal assurance remain open.

\section{A concrete cutoff for each multiplier}
\begin{theorem}\label{thm:cutoff}
For $a\ge15$ put
\[
 P_0(a)=\left\lceil
 \frac{\max\{400000,\exp\sqrt{a/50}\}}{a-2}\right\rceil. \tag{1}
\]
For every prime $p>6a$, $r\ge1$, $P=p^r\ge P_0(a)$ and
$0\le i\le12$, the polynomial $F_{aP+i}$ is squarefree.
For $a\le14$ the retained uniform family already applies without a
cutoff under its condition $p>6\max(a,i)$.
\end{theorem}
\begin{proof}
Put $n=aP+i$ and $x=n-2P=(a-2)P+i$.
\href{sources/dusart-1002.0442.pdf}{Dusart, Proposition 6.8}, supplies
an unconditional prime
\[
 x<q\le x+\frac{x}{25\log^2x}\quad(x\ge396738).
\]
The primary arXiv PDF and its extracted text are pinned; this published
estimate is imported rather than its numerical proof reimplemented.
By (1), $x\ge400000$ and $\log^2x\ge a/50$. Also $i<2P$ because
$p>6a\ge90$ and $i\le12$. Hence $x<aP$ and
\[
 \frac{x}{25\log^2x}\le\frac{2x}{a}<2P,
 \qquad n-2P<q<n. \tag{2}
\]
For the largest prime $q_n\le n$ we have $q_n\ge q$.
The retained prime-gap theorem therefore bounds every rational square
divisor $H$ by $\deg H\le(n-q_n)/2<P$.
The original shifted $p$-polygon forces $\deg H=h+Pv$.
Thus $v=0$. Its horizontal residual is $F_i\bmod p$; the retained
certificate at $i\le12$, $p>6a>6i$, gives $h\le1$.
Real-place parity excludes one, so no repeated rational irreducible
factor exists. Since $F_n$ is rational, it is squarefree over $\C$ as well.
\end{proof}

Exactly,
\[
 P_0(15)=30770,\qquad P_0(100)=4082,\qquad P_0(1000)=401. \tag{3}
\]
For $a=1000$, $p>6000$ already exceeds the cutoff, so the whole
prime-conditioned family is covered. For any fixed $a$, potential
exceptions to this cutoff criterion lie in the explicit finite set
$p^r<P_0(a)$, $p>6a$. These are inputs unclassified by this cutoff
criterion alone, not counterexamples. Corollary~\ref{cor:bounded-multipliers}
below discharges all of them when $a\le24000$.

The exporter evaluates (1) without floating-point rounding assumptions.
It compares $50\log^2((a-2)m)$ to $a$ using exact rational brackets
from $\log y=2\sum_{j\ge0}t^{2j+1}/(2j+1)$,
$t=(y-1)/(y+1)$, after binary range reduction $1\le y<2$.
The positive tail after $N$ terms is at most
$2t^{2N+1}/((2N+1)(1-t^2))$. An integer upper bracket uses
$e^{\sqrt{a/50}}<3^{\lceil\sqrt{a/50}\rceil}$.
Binary search and certified comparisons give the exact ceiling.
Equality cannot occur at an integer comparison point: the exponential
of a nonzero algebraic number is transcendental, as already used in the
\href{COUPLED_BLOCK_PERIODS.tex}{retained exponential-period proof}.
The packet retains the final true/false rational brackets.

\section{All initial exceptions through multiplier 24000}
\begin{corollary}\label{cor:bounded-multipliers}
For every integer $1\le a\le24000$, every $0\le i\le12$, every prime
$p>6\max(a,i)$ and every $r\ge1$, the polynomial $F_{ap^r+i}$ is
squarefree. There are no excluded initial prime powers in this range.
\end{corollary}
\begin{proof}
The retained uniform family handles $a\le14$. Assume $15\le a\le24000$,
put $P=p^r$, $n=aP+i$ and $x=n-2P$. Since $P>6a$,
$x>6a(a-2)$. For real $a\ge3$ write
\[
 L(a)=\log(6a(a-2)),\quad H(a)=50L(a)^2-a,\quad
 S=a^{-2}+(a-2)^{-2}.
\]
Then $L''=-S$, $(L')^2\le2S$, and $L(a)>2$, giving
\[
 H''(a)=100((L')^2-LS)\le100(2-L)S<0.
\]
Hence $H$ is strictly concave and lies above its endpoint chord.
The following are exact strict bounds:
\[
 L(15)>6,\quad H(15)>1785,\qquad
 L(24000)>\frac{263}{12},\quad H(24000)>\frac{1225}{72}>0.
\]
For completeness, the positive exponential series gives
\[
 e<\sum_{j=0}^{8}\frac1{j!}+\frac{10}{9\cdot9!}
   <\frac{2719}{1000}<3.
\]
The geometric remainder bound uses consecutive omitted-term ratios
at most $1/10$. Now $e^6<3^6<1170=6\cdot15\cdot13$.
Also $e^{1/12}>1+1/12=13/12$, so the retained rational integer comparison
\[
 e^{263/12}=\frac{e^{22}}{e^{1/12}}
 <\left(\frac{2719}{1000}\right)^{22}\frac{12}{13}
 <3455712000=6\cdot24000\cdot23998
\]
proves the right endpoint bound. Concavity therefore gives $H(a)>0$
throughout the entire real interval $[15,24000]$. As $x>6a(a-2)>1$,
\[
 \log^2x>a/50.
\]
If $x\ge400000$, the same Dusart argument as in Theorem~\ref{thm:cutoff}
gives $n-2P<q<n$ and thus $\deg G<P$ for any rational $G^2\mid F_n$.
This branch no longer requires $P\ge P_0(a)$.

If $x<400000$, then
\[
 n=x+2P\le\frac{a}{a-2}x
 <400000\frac{15}{13}<461539.
\]
The complete prime list through $500000$ is retained in
\href{initial_prime_gaps_500000.json}{the finite gap certificate}.
An Eratosthenes sieve and a separate Euler least-prime-factor sieve
agree on every prime. They give $41538$ primes, maximum consecutive
prime gap $114$, attained at $(492113,492227)$, and
\[
 \max_{2\le m\le500000}(m-q_m)=113.
\]
The last-prime tail is included; no unrecorded prime beyond the upper
limit is needed. Here $n\ge15\cdot97>30$ and $P>90$, so
$n-q_n\le113<2P$. The prime-gap theorem again gives $\deg G<P$.
The finite computation is a prime-gap certificate, not a scan of
factorial moment polynomials or an all-index squarefreeness assertion.

In both branches the shifted $p$-polygon forces the sloped contribution
in $\deg G=h+Pv$ to vanish. The retained horizontal kernel certificate
for $i\le12$ gives $h\le1$, and real-place parity excludes degree one.
Thus $F_n$ is squarefree in every case claimed.
\end{proof}
The exact endpoint comparisons are retained in
\href{multiplier_endpoint_certificate.json}{a rational certificate}.
This corollary discharges all initial cutoff exceptions for $a\le24000$.
Unbounded multipliers and unbounded offsets remain outside its scope;
the existing explicit cutoff still applies for larger fixed multipliers.

\section{What the unbounded integer-kernel target would prove}
For $d\ge2$ retain
\[
 R_d=\operatorname{Res}(F_d,F_d'),\quad
 \operatorname{Sres}_1(F_d,F_d')=A_dT+B_d,\quad
 C_d=\gcd(R_d,A_d,B_d)\ge0.
\]
\begin{lemma}\label{lem:strength}
For this rational moment sequence,
\[
 C_d>0\quad\Longleftrightarrow\quad F_d\text{ is squarefree over }\Q.
 \tag{4}
\]
Consequently, proving positivity for every $d$ already proves the
all-index characteristic-zero squarefreeness statement. Proving in
addition that all prime divisors of every $C_d$ are at most $6d$ imposes
the stronger modular-collision restriction from the retained criterion.
\end{lemma}
\begin{proof}
If $F_d$ is squarefree then $R_d\ne0$, so $C_d>0$.
If it is not squarefree, a repeated rational irreducible factor has
even degree at least two by the real-place theorem. Thus
$\gcd(F_d,F_d')$ has degree at least two. The resultant and the
first subresultant both vanish, so $R_d=A_d=B_d=0$ and $C_d=0$.
\end{proof}
The finite list through twelve generates infinite prime-conditioned
families without asserting the unbounded statement (4).

\section{Full residual square-part degrees and the weaker coverage target}
Let $n=q+d$, where $q>6\max(1,d)$ is prime. The retained polygon
has a horizontal edge of width $d$ and a primitive width-$q$ edge.
For any rational square divisor, reduction preserves its degree on
that horizontal edge and square-divides $\overline{F_d}$.
Factor
\[
 \overline{F_d}=u\prod_\phi\phi^{e_\phi},\qquad
 K_{d,q}=\prod_\phi\phi^{\lfloor e_\phi/2\rfloor},
 \qquad
 D_{d,q}=\left\{\sum_\phi\deg(\phi)t_\phi:
 0\le t_\phi\le\lfloor e_\phi/2\rfloor\right\}. \tag{5}
\]
All factors are monic irreducibles over $\F_q$. The scalar $u$ is irrelevant.

\begin{theorem}\label{thm:degrees}
For any finite nonempty set $S$ of primes $6n/7<q\le n$, $q>6$,
a rational square divisor of $F_n$ has degree in
\[
 \bigcap_{q\in S}D_{n-q,q}.
\]
Intersect this with the nonnegative even degrees and every valid global
bound, including $2\lfloor(n-q_n)/4\rfloor$ for $n\ge30$.
If only zero remains, $F_n$ is squarefree. A surviving degree makes
no counterexample claim.
\end{theorem}
\begin{proof}
At each prime, doubled weighted initial multiplication makes the
reduced horizontal divisor divide $K_{n-q,q}$ with its degree retained.
The possible divisor degrees of $K$ are exactly the bounded sum (5).
The same rational divisor has one common degree at all primes.
Its degree is even and obeys the retained global bounds, proving every
necessary restriction and the asserted sufficient zero-only test.
\end{proof}
For example, a square part that is one irreducible cubic permits only
$0,3$. Parity excludes its positive degree even though the residual
has a substantial repeated factor. This illustrates the degree criterion;
it is not a claim that a synthetic residual is an actual moment kernel.

The scalar shortcut $q\nmid C_{n-q}$ implies residual derivative-gcd
degree at most one. Offsets zero and one are automatically good.
Thus proving, for every $n\ge N$, the existence of a nearby prime
$6n/7<q\le n$ with that scalar property would prove squarefreeness
for all $n\ge N$. Full all-index completion would additionally need
the finite initial indices. More permissively, it suffices that the full
degree intersections above exclude every positive even degree; kernels
may have exceptional large primes and may have derivative gcd degree
larger than one.

For a fixed prime $p>6\max(1,i)$ the same residual degree test on
$F_i\bmod p$ also applies to $F_{p^r+i}$ for every $r\ge1$.
This can certify individual fixed-prime offset families beyond the
retained uniform prefix, without proving smoothness of $C_i$ at all primes.
For example, the retained exact factorization of $F_{13}\bmod113$ has
two distinct irreducible factors of degrees one and twelve, so its square
part is one. Reconstruction and Frobenius irreducibility witnesses
therefore certify every $F_{113^r+13}$, $r\ge1$. This fixed-prime
example makes no all-prime assertion about $C_{13}$.

\section{Executable boundary}
\href{effective_residual_windows.py}{The exporter} retains exact cutoff
comparisons, the complete multiplier range through 24000, finite exception
bounds for larger multipliers, full residual factorizations and
square-part divisor degrees. It never enumerates a large nearby-prime
range implicitly: finite prime inputs are explicit.
\href{verify_effective_residual_windows.py}{Controls} replay rational
logarithm brackets, strict cutoff inequalities, independent modular moments,
factor reconstruction and Frobenius irreducibility checks. Integer
thresholds are not inferred from rounded transcendental numerics.

The fixed-multiplier effectiveness gap is closed. The substantive remaining
all-index step is uniform coverage by the residual degree intersections,
with the finite initial range handled as part of any complete proof.
The unbounded $C_d$ positivity or smoothness assertions are not reported as
proved, and there is no Gamma-shape or novelty claim.
\end{document}
